% ID: SEL-D0026-OA-E02
% ESTADO: validada. S027, cotejada con examen y criterios.
\textbf{a)} Los tramos son polinómicos. En $2$, para que el valor
$4/a+1$ coincida con el límite derecho $a-2$, debe cumplirse
$a^2-3a-4=(a-4)(a+1)=0$. Como $a>0$, $\boxed{a=4}$.
Las pendientes laterales en $2$ son $4/a=1$ y $-1$.
La función es continua en todo $\mathbb R$ y derivable salvo en $2$.

\textbf{b)} La parábola $x^2/4+1$ tiene vértice $(0,1)$
y enlaza con la recta $y=4-x$ en $(2,2)$.
En $x=-1$, el punto es $(-1,5/4)$ y la pendiente es $-1/2$.
Así, $y-5/4=-(x+1)/2$, o $\boxed{y=3/4-x/2}$.
\begin{center}
\setlength{\unitlength}{8mm}
\begin{picture}(14,8)(-7,-2)
\qbezier(-4.5,0)(0.5,0)(5.5,0)\put(5.6,0){$x$}
\qbezier(0,-1.5)(0,2)(0,5.5)\put(0,5.6){$y$}
\qbezier(-4,5)(-1,-1)(2,2)
\qbezier(2,2)(3.5,0.5)(5,-1)
\qbezier(-4,2.75)(0,0.75)(4,-1.25)
\put(-4.4,3){\makebox(0,0)[r]{\small tangente}}
\put(2.2,2.2){$(2,2)$}
\put(-1,1.35){\makebox(0,0)[b]{$(-1,5/4)$}}
\end{picture}
\end{center}
